% ─────────────────────────────────────────────────────────────────────────────
% SHARED coursework preamble — every course inputs THIS.
%
% PREAMBLE ONLY. No \begin{document}, no course-specific notation.
%
% Course-specific macros live at COURSE level, in
%   content/<course>/reference_docs/<course>_macros.tex
% which is inputted AFTER this file. ONE per course, shared by hw/qz/exam --
% macros under a single kind are unreachable from a sibling kind.
%
% That split is the point: the page style, the answer box, and the problem
% header are identical everywhere, so they are defined once; \conv and \ustep
% mean something in DSP and nothing in a networks course, so they are not.
%
% A per-problem file therefore starts:
%   % ─────────────────────────────────────────────────────────────────────────────
% SHARED coursework preamble — every course inputs THIS.
%
% PREAMBLE ONLY. No \begin{document}, no course-specific notation.
%
% Course-specific macros live at COURSE level, in
%   content/<course>/reference_docs/<course>_macros.tex
% which is inputted AFTER this file. ONE per course, shared by hw/qz/exam --
% macros under a single kind are unreachable from a sibling kind.
%
% That split is the point: the page style, the answer box, and the problem
% header are identical everywhere, so they are defined once; \conv and \ustep
% mean something in DSP and nothing in a networks course, so they are not.
%
% A per-problem file therefore starts:
%   % ─────────────────────────────────────────────────────────────────────────────
% SHARED coursework preamble — every course inputs THIS.
%
% PREAMBLE ONLY. No \begin{document}, no course-specific notation.
%
% Course-specific macros live at COURSE level, in
%   content/<course>/reference_docs/<course>_macros.tex
% which is inputted AFTER this file. ONE per course, shared by hw/qz/exam --
% macros under a single kind are unreachable from a sibling kind.
%
% That split is the point: the page style, the answer box, and the problem
% header are identical everywhere, so they are defined once; \conv and \ustep
% mean something in DSP and nothing in a networks course, so they are not.
%
% A per-problem file therefore starts:
%   \input{../../../../docs/latex/coursework_preamble.tex}
%   \input{../../reference_docs/cesc470_macros.tex}
%   \begin{document}
%
% Build-check one file, or a whole assignment:
%   docs/latex/build_tex.sh content/cesc_470/hw/hw01/p01_five_components.tex
%   docs/latex/build_tex.sh content/cesc_470/hw/hw01
%
% Doctrine: docs/directives/coursework-solutions.md
% ─────────────────────────────────────────────────────────────────────────────

\documentclass[11pt]{article}

\usepackage[margin=1in]{geometry}
\usepackage{amsmath, amssymb, mathtools}
\usepackage{graphicx}
\usepackage{float}
\usepackage{booktabs}
\usepackage{longtable}
\usepackage{enumitem}
\usepackage{siunitx}
\usepackage[dvipsnames]{xcolor}
\usepackage{tcolorbox}
\usepackage{fancyhdr}
\usepackage{caption}
\usepackage{tikz}
\usepackage{pgfplots}
\pgfplotsset{compat=1.18}
\usepackage[colorlinks=true, allcolors=NavyBlue]{hyperref}

\captionsetup{font=small, labelfont=bf}
\setlist[enumerate]{itemsep=2pt, topsep=4pt}
\setlist[itemize]{itemsep=2pt, topsep=4pt}

% --- final-answer box -------------------------------------------------------
% Every problem file ends with its result in one of these. The condensed
% solutions document is assembled by lifting exactly these.
\newtcolorbox{answerbox}[1][]{
  colback=Green!6, colframe=Green!55!black,
  boxrule=0.6pt, arc=2pt, left=6pt, right=6pt, top=4pt, bottom=4pt,
  #1
}

% --- citation to course material --------------------------------------------
% \cite{Module 01, PDF p55}  ->  a small italic parenthetical.
% Solutions cite the SLIDE/PAGE a formula came from so a grader (or a future
% reader) can check the claim against what was actually taught. See the
% directive: every non-obvious step carries one.
\newcommand{\srcref}[1]{{\small\itshape(#1)}}

% --- a step that came from OUTSIDE the course materials ---------------------
% Used only where the course is genuinely silent and the problem asks for
% outside research. Marked visually so it is never mistaken for lecture content.
\newcommand{\extref}[1]{{\small\itshape\color{Sepia}[external: #1]}}

% --- callout for the trap in a problem --------------------------------------
\newtcolorbox{trap}[1][]{
  colback=BurntOrange!7, colframe=BurntOrange!70!black,
  boxrule=0.6pt, arc=2pt, left=6pt, right=6pt, top=4pt, bottom=4pt,
  fonttitle=\bfseries, title=Watch out, #1
}

% --- per-problem header -----------------------------------------------------
% \problemheader{8}{15 pts}{Target clock rate}
% The optional 4-argument form carries a learning outcome where the course
% states one (CESC 410 does; CESC 470's HW1 does not).
\newcommand{\problemheader}[3]{%
  \begin{center}\large\bfseries
    Problem #1 \quad\normalsize\mdseries(#2)\\[2pt]
    \large\bfseries #3
  \end{center}
  \vspace{2pt}\hrule\vspace{8pt}
}
\newcommand{\problemheaderlo}[4]{%
  \begin{center}\large\bfseries
    Problem #1 \quad\normalsize\mdseries(#2, #3)\\[2pt]
    \large\bfseries #4
  \end{center}
  \vspace{2pt}\hrule\vspace{8pt}
}

% Course name in the running head. Defined BEFORE \lhead references it, and
% \renewcommand'd by each course's macros file (inputted after this one).
\providecommand{\coursename}{}

\pagestyle{fancy}
\fancyhf{}
\lhead{\small\coursename}
\rhead{\small\thepage}
\renewcommand{\headrulewidth}{0.3pt}

%   % CESC 470 Computer Architecture — course-specific macros ONLY.
%
% Inputted AFTER docs/latex/coursework_preamble.tex, which carries everything
% shared (answerbox, problemheader, srcref, page style). Put nothing general
% here: if another course would want it, it belongs in the shared preamble.

\renewcommand{\coursename}{CESC 470}

% --- performance-equation notation ------------------------------------------
% The course writes these in words on the slides; these render them compactly
% and, more importantly, IDENTICALLY across every problem file.
\newcommand{\CPUtime}{T_{\mathrm{CPU}}}
\newcommand{\IC}{\mathrm{IC}}                    % instruction count
\newcommand{\CPI}{\mathrm{CPI}}
\newcommand{\cycletime}{t_{\mathrm{cycle}}}
\newcommand{\clockrate}{f_{\mathrm{clk}}}
\newcommand{\cycles}{N_{\mathrm{cycles}}}
\newcommand{\speedup}{S}

% Amdahl's law, in the form the slides state it (Module 01, PDF p86, slide 49).
\newcommand{\amdahl}[3]{#1 + \dfrac{#2}{#3}}

%   \begin{document}
%
% Build-check one file, or a whole assignment:
%   docs/latex/build_tex.sh content/cesc_470/hw/hw01/p01_five_components.tex
%   docs/latex/build_tex.sh content/cesc_470/hw/hw01
%
% Doctrine: docs/directives/coursework-solutions.md
% ─────────────────────────────────────────────────────────────────────────────

\documentclass[11pt]{article}

\usepackage[margin=1in]{geometry}
\usepackage{amsmath, amssymb, mathtools}
\usepackage{graphicx}
\usepackage{float}
\usepackage{booktabs}
\usepackage{longtable}
\usepackage{enumitem}
\usepackage{siunitx}
\usepackage[dvipsnames]{xcolor}
\usepackage{tcolorbox}
\usepackage{fancyhdr}
\usepackage{caption}
\usepackage{tikz}
\usepackage{pgfplots}
\pgfplotsset{compat=1.18}
\usepackage[colorlinks=true, allcolors=NavyBlue]{hyperref}

\captionsetup{font=small, labelfont=bf}
\setlist[enumerate]{itemsep=2pt, topsep=4pt}
\setlist[itemize]{itemsep=2pt, topsep=4pt}

% --- final-answer box -------------------------------------------------------
% Every problem file ends with its result in one of these. The condensed
% solutions document is assembled by lifting exactly these.
\newtcolorbox{answerbox}[1][]{
  colback=Green!6, colframe=Green!55!black,
  boxrule=0.6pt, arc=2pt, left=6pt, right=6pt, top=4pt, bottom=4pt,
  #1
}

% --- citation to course material --------------------------------------------
% \cite{Module 01, PDF p55}  ->  a small italic parenthetical.
% Solutions cite the SLIDE/PAGE a formula came from so a grader (or a future
% reader) can check the claim against what was actually taught. See the
% directive: every non-obvious step carries one.
\newcommand{\srcref}[1]{{\small\itshape(#1)}}

% --- a step that came from OUTSIDE the course materials ---------------------
% Used only where the course is genuinely silent and the problem asks for
% outside research. Marked visually so it is never mistaken for lecture content.
\newcommand{\extref}[1]{{\small\itshape\color{Sepia}[external: #1]}}

% --- callout for the trap in a problem --------------------------------------
\newtcolorbox{trap}[1][]{
  colback=BurntOrange!7, colframe=BurntOrange!70!black,
  boxrule=0.6pt, arc=2pt, left=6pt, right=6pt, top=4pt, bottom=4pt,
  fonttitle=\bfseries, title=Watch out, #1
}

% --- per-problem header -----------------------------------------------------
% \problemheader{8}{15 pts}{Target clock rate}
% The optional 4-argument form carries a learning outcome where the course
% states one (CESC 410 does; CESC 470's HW1 does not).
\newcommand{\problemheader}[3]{%
  \begin{center}\large\bfseries
    Problem #1 \quad\normalsize\mdseries(#2)\\[2pt]
    \large\bfseries #3
  \end{center}
  \vspace{2pt}\hrule\vspace{8pt}
}
\newcommand{\problemheaderlo}[4]{%
  \begin{center}\large\bfseries
    Problem #1 \quad\normalsize\mdseries(#2, #3)\\[2pt]
    \large\bfseries #4
  \end{center}
  \vspace{2pt}\hrule\vspace{8pt}
}

% Course name in the running head. Defined BEFORE \lhead references it, and
% \renewcommand'd by each course's macros file (inputted after this one).
\providecommand{\coursename}{}

\pagestyle{fancy}
\fancyhf{}
\lhead{\small\coursename}
\rhead{\small\thepage}
\renewcommand{\headrulewidth}{0.3pt}

%   % CESC 470 Computer Architecture — course-specific macros ONLY.
%
% Inputted AFTER docs/latex/coursework_preamble.tex, which carries everything
% shared (answerbox, problemheader, srcref, page style). Put nothing general
% here: if another course would want it, it belongs in the shared preamble.

\renewcommand{\coursename}{CESC 470}

% --- performance-equation notation ------------------------------------------
% The course writes these in words on the slides; these render them compactly
% and, more importantly, IDENTICALLY across every problem file.
\newcommand{\CPUtime}{T_{\mathrm{CPU}}}
\newcommand{\IC}{\mathrm{IC}}                    % instruction count
\newcommand{\CPI}{\mathrm{CPI}}
\newcommand{\cycletime}{t_{\mathrm{cycle}}}
\newcommand{\clockrate}{f_{\mathrm{clk}}}
\newcommand{\cycles}{N_{\mathrm{cycles}}}
\newcommand{\speedup}{S}

% Amdahl's law, in the form the slides state it (Module 01, PDF p86, slide 49).
\newcommand{\amdahl}[3]{#1 + \dfrac{#2}{#3}}

%   \begin{document}
%
% Build-check one file, or a whole assignment:
%   docs/latex/build_tex.sh content/cesc_470/hw/hw01/p01_five_components.tex
%   docs/latex/build_tex.sh content/cesc_470/hw/hw01
%
% Doctrine: docs/directives/coursework-solutions.md
% ─────────────────────────────────────────────────────────────────────────────

\documentclass[11pt]{article}

\usepackage[margin=1in]{geometry}
\usepackage{amsmath, amssymb, mathtools}
\usepackage{graphicx}
\usepackage{float}
\usepackage{booktabs}
\usepackage{longtable}
\usepackage{enumitem}
\usepackage{siunitx}
\usepackage[dvipsnames]{xcolor}
\usepackage{tcolorbox}
\usepackage{fancyhdr}
\usepackage{caption}
\usepackage{tikz}
\usepackage{pgfplots}
\pgfplotsset{compat=1.18}
\usepackage[colorlinks=true, allcolors=NavyBlue]{hyperref}

\captionsetup{font=small, labelfont=bf}
\setlist[enumerate]{itemsep=2pt, topsep=4pt}
\setlist[itemize]{itemsep=2pt, topsep=4pt}

% --- final-answer box -------------------------------------------------------
% Every problem file ends with its result in one of these. The condensed
% solutions document is assembled by lifting exactly these.
\newtcolorbox{answerbox}[1][]{
  colback=Green!6, colframe=Green!55!black,
  boxrule=0.6pt, arc=2pt, left=6pt, right=6pt, top=4pt, bottom=4pt,
  #1
}

% --- citation to course material --------------------------------------------
% \cite{Module 01, PDF p55}  ->  a small italic parenthetical.
% Solutions cite the SLIDE/PAGE a formula came from so a grader (or a future
% reader) can check the claim against what was actually taught. See the
% directive: every non-obvious step carries one.
\newcommand{\srcref}[1]{{\small\itshape(#1)}}

% --- a step that came from OUTSIDE the course materials ---------------------
% Used only where the course is genuinely silent and the problem asks for
% outside research. Marked visually so it is never mistaken for lecture content.
\newcommand{\extref}[1]{{\small\itshape\color{Sepia}[external: #1]}}

% --- callout for the trap in a problem --------------------------------------
\newtcolorbox{trap}[1][]{
  colback=BurntOrange!7, colframe=BurntOrange!70!black,
  boxrule=0.6pt, arc=2pt, left=6pt, right=6pt, top=4pt, bottom=4pt,
  fonttitle=\bfseries, title=Watch out, #1
}

% --- per-problem header -----------------------------------------------------
% \problemheader{8}{15 pts}{Target clock rate}
% The optional 4-argument form carries a learning outcome where the course
% states one (CESC 410 does; CESC 470's HW1 does not).
\newcommand{\problemheader}[3]{%
  \begin{center}\large\bfseries
    Problem #1 \quad\normalsize\mdseries(#2)\\[2pt]
    \large\bfseries #3
  \end{center}
  \vspace{2pt}\hrule\vspace{8pt}
}
\newcommand{\problemheaderlo}[4]{%
  \begin{center}\large\bfseries
    Problem #1 \quad\normalsize\mdseries(#2, #3)\\[2pt]
    \large\bfseries #4
  \end{center}
  \vspace{2pt}\hrule\vspace{8pt}
}

% Course name in the running head. Defined BEFORE \lhead references it, and
% \renewcommand'd by each course's macros file (inputted after this one).
\providecommand{\coursename}{}

\pagestyle{fancy}
\fancyhf{}
\lhead{\small\coursename}
\rhead{\small\thepage}
\renewcommand{\headrulewidth}{0.3pt}

% ─────────────────────────────────────────────────────────────────────────────
% CESC 410 course-specific macros. Inputted AFTER the shared preamble:
%
%   % ─────────────────────────────────────────────────────────────────────────────
% SHARED coursework preamble — every course inputs THIS.
%
% PREAMBLE ONLY. No \begin{document}, no course-specific notation.
%
% Course-specific macros live at COURSE level, in
%   content/<course>/reference_docs/<course>_macros.tex
% which is inputted AFTER this file. ONE per course, shared by hw/qz/exam --
% macros under a single kind are unreachable from a sibling kind.
%
% That split is the point: the page style, the answer box, and the problem
% header are identical everywhere, so they are defined once; \conv and \ustep
% mean something in DSP and nothing in a networks course, so they are not.
%
% A per-problem file therefore starts:
%   % ─────────────────────────────────────────────────────────────────────────────
% SHARED coursework preamble — every course inputs THIS.
%
% PREAMBLE ONLY. No \begin{document}, no course-specific notation.
%
% Course-specific macros live at COURSE level, in
%   content/<course>/reference_docs/<course>_macros.tex
% which is inputted AFTER this file. ONE per course, shared by hw/qz/exam --
% macros under a single kind are unreachable from a sibling kind.
%
% That split is the point: the page style, the answer box, and the problem
% header are identical everywhere, so they are defined once; \conv and \ustep
% mean something in DSP and nothing in a networks course, so they are not.
%
% A per-problem file therefore starts:
%   \input{../../../../docs/latex/coursework_preamble.tex}
%   \input{../../reference_docs/cesc470_macros.tex}
%   \begin{document}
%
% Build-check one file, or a whole assignment:
%   docs/latex/build_tex.sh content/cesc_470/hw/hw01/p01_five_components.tex
%   docs/latex/build_tex.sh content/cesc_470/hw/hw01
%
% Doctrine: docs/directives/coursework-solutions.md
% ─────────────────────────────────────────────────────────────────────────────

\documentclass[11pt]{article}

\usepackage[margin=1in]{geometry}
\usepackage{amsmath, amssymb, mathtools}
\usepackage{graphicx}
\usepackage{float}
\usepackage{booktabs}
\usepackage{longtable}
\usepackage{enumitem}
\usepackage{siunitx}
\usepackage[dvipsnames]{xcolor}
\usepackage{tcolorbox}
\usepackage{fancyhdr}
\usepackage{caption}
\usepackage{tikz}
\usepackage{pgfplots}
\pgfplotsset{compat=1.18}
\usepackage[colorlinks=true, allcolors=NavyBlue]{hyperref}

\captionsetup{font=small, labelfont=bf}
\setlist[enumerate]{itemsep=2pt, topsep=4pt}
\setlist[itemize]{itemsep=2pt, topsep=4pt}

% --- final-answer box -------------------------------------------------------
% Every problem file ends with its result in one of these. The condensed
% solutions document is assembled by lifting exactly these.
\newtcolorbox{answerbox}[1][]{
  colback=Green!6, colframe=Green!55!black,
  boxrule=0.6pt, arc=2pt, left=6pt, right=6pt, top=4pt, bottom=4pt,
  #1
}

% --- citation to course material --------------------------------------------
% \cite{Module 01, PDF p55}  ->  a small italic parenthetical.
% Solutions cite the SLIDE/PAGE a formula came from so a grader (or a future
% reader) can check the claim against what was actually taught. See the
% directive: every non-obvious step carries one.
\newcommand{\srcref}[1]{{\small\itshape(#1)}}

% --- a step that came from OUTSIDE the course materials ---------------------
% Used only where the course is genuinely silent and the problem asks for
% outside research. Marked visually so it is never mistaken for lecture content.
\newcommand{\extref}[1]{{\small\itshape\color{Sepia}[external: #1]}}

% --- callout for the trap in a problem --------------------------------------
\newtcolorbox{trap}[1][]{
  colback=BurntOrange!7, colframe=BurntOrange!70!black,
  boxrule=0.6pt, arc=2pt, left=6pt, right=6pt, top=4pt, bottom=4pt,
  fonttitle=\bfseries, title=Watch out, #1
}

% --- per-problem header -----------------------------------------------------
% \problemheader{8}{15 pts}{Target clock rate}
% The optional 4-argument form carries a learning outcome where the course
% states one (CESC 410 does; CESC 470's HW1 does not).
\newcommand{\problemheader}[3]{%
  \begin{center}\large\bfseries
    Problem #1 \quad\normalsize\mdseries(#2)\\[2pt]
    \large\bfseries #3
  \end{center}
  \vspace{2pt}\hrule\vspace{8pt}
}
\newcommand{\problemheaderlo}[4]{%
  \begin{center}\large\bfseries
    Problem #1 \quad\normalsize\mdseries(#2, #3)\\[2pt]
    \large\bfseries #4
  \end{center}
  \vspace{2pt}\hrule\vspace{8pt}
}

% Course name in the running head. Defined BEFORE \lhead references it, and
% \renewcommand'd by each course's macros file (inputted after this one).
\providecommand{\coursename}{}

\pagestyle{fancy}
\fancyhf{}
\lhead{\small\coursename}
\rhead{\small\thepage}
\renewcommand{\headrulewidth}{0.3pt}

%   % CESC 470 Computer Architecture — course-specific macros ONLY.
%
% Inputted AFTER docs/latex/coursework_preamble.tex, which carries everything
% shared (answerbox, problemheader, srcref, page style). Put nothing general
% here: if another course would want it, it belongs in the shared preamble.

\renewcommand{\coursename}{CESC 470}

% --- performance-equation notation ------------------------------------------
% The course writes these in words on the slides; these render them compactly
% and, more importantly, IDENTICALLY across every problem file.
\newcommand{\CPUtime}{T_{\mathrm{CPU}}}
\newcommand{\IC}{\mathrm{IC}}                    % instruction count
\newcommand{\CPI}{\mathrm{CPI}}
\newcommand{\cycletime}{t_{\mathrm{cycle}}}
\newcommand{\clockrate}{f_{\mathrm{clk}}}
\newcommand{\cycles}{N_{\mathrm{cycles}}}
\newcommand{\speedup}{S}

% Amdahl's law, in the form the slides state it (Module 01, PDF p86, slide 49).
\newcommand{\amdahl}[3]{#1 + \dfrac{#2}{#3}}

%   \begin{document}
%
% Build-check one file, or a whole assignment:
%   docs/latex/build_tex.sh content/cesc_470/hw/hw01/p01_five_components.tex
%   docs/latex/build_tex.sh content/cesc_470/hw/hw01
%
% Doctrine: docs/directives/coursework-solutions.md
% ─────────────────────────────────────────────────────────────────────────────

\documentclass[11pt]{article}

\usepackage[margin=1in]{geometry}
\usepackage{amsmath, amssymb, mathtools}
\usepackage{graphicx}
\usepackage{float}
\usepackage{booktabs}
\usepackage{longtable}
\usepackage{enumitem}
\usepackage{siunitx}
\usepackage[dvipsnames]{xcolor}
\usepackage{tcolorbox}
\usepackage{fancyhdr}
\usepackage{caption}
\usepackage{tikz}
\usepackage{pgfplots}
\pgfplotsset{compat=1.18}
\usepackage[colorlinks=true, allcolors=NavyBlue]{hyperref}

\captionsetup{font=small, labelfont=bf}
\setlist[enumerate]{itemsep=2pt, topsep=4pt}
\setlist[itemize]{itemsep=2pt, topsep=4pt}

% --- final-answer box -------------------------------------------------------
% Every problem file ends with its result in one of these. The condensed
% solutions document is assembled by lifting exactly these.
\newtcolorbox{answerbox}[1][]{
  colback=Green!6, colframe=Green!55!black,
  boxrule=0.6pt, arc=2pt, left=6pt, right=6pt, top=4pt, bottom=4pt,
  #1
}

% --- citation to course material --------------------------------------------
% \cite{Module 01, PDF p55}  ->  a small italic parenthetical.
% Solutions cite the SLIDE/PAGE a formula came from so a grader (or a future
% reader) can check the claim against what was actually taught. See the
% directive: every non-obvious step carries one.
\newcommand{\srcref}[1]{{\small\itshape(#1)}}

% --- a step that came from OUTSIDE the course materials ---------------------
% Used only where the course is genuinely silent and the problem asks for
% outside research. Marked visually so it is never mistaken for lecture content.
\newcommand{\extref}[1]{{\small\itshape\color{Sepia}[external: #1]}}

% --- callout for the trap in a problem --------------------------------------
\newtcolorbox{trap}[1][]{
  colback=BurntOrange!7, colframe=BurntOrange!70!black,
  boxrule=0.6pt, arc=2pt, left=6pt, right=6pt, top=4pt, bottom=4pt,
  fonttitle=\bfseries, title=Watch out, #1
}

% --- per-problem header -----------------------------------------------------
% \problemheader{8}{15 pts}{Target clock rate}
% The optional 4-argument form carries a learning outcome where the course
% states one (CESC 410 does; CESC 470's HW1 does not).
\newcommand{\problemheader}[3]{%
  \begin{center}\large\bfseries
    Problem #1 \quad\normalsize\mdseries(#2)\\[2pt]
    \large\bfseries #3
  \end{center}
  \vspace{2pt}\hrule\vspace{8pt}
}
\newcommand{\problemheaderlo}[4]{%
  \begin{center}\large\bfseries
    Problem #1 \quad\normalsize\mdseries(#2, #3)\\[2pt]
    \large\bfseries #4
  \end{center}
  \vspace{2pt}\hrule\vspace{8pt}
}

% Course name in the running head. Defined BEFORE \lhead references it, and
% \renewcommand'd by each course's macros file (inputted after this one).
\providecommand{\coursename}{}

\pagestyle{fancy}
\fancyhf{}
\lhead{\small\coursename}
\rhead{\small\thepage}
\renewcommand{\headrulewidth}{0.3pt}

%   % ─────────────────────────────────────────────────────────────────────────────
% CESC 410 course-specific macros. Inputted AFTER the shared preamble:
%
%   % ─────────────────────────────────────────────────────────────────────────────
% SHARED coursework preamble — every course inputs THIS.
%
% PREAMBLE ONLY. No \begin{document}, no course-specific notation.
%
% Course-specific macros live at COURSE level, in
%   content/<course>/reference_docs/<course>_macros.tex
% which is inputted AFTER this file. ONE per course, shared by hw/qz/exam --
% macros under a single kind are unreachable from a sibling kind.
%
% That split is the point: the page style, the answer box, and the problem
% header are identical everywhere, so they are defined once; \conv and \ustep
% mean something in DSP and nothing in a networks course, so they are not.
%
% A per-problem file therefore starts:
%   \input{../../../../docs/latex/coursework_preamble.tex}
%   \input{../../reference_docs/cesc470_macros.tex}
%   \begin{document}
%
% Build-check one file, or a whole assignment:
%   docs/latex/build_tex.sh content/cesc_470/hw/hw01/p01_five_components.tex
%   docs/latex/build_tex.sh content/cesc_470/hw/hw01
%
% Doctrine: docs/directives/coursework-solutions.md
% ─────────────────────────────────────────────────────────────────────────────

\documentclass[11pt]{article}

\usepackage[margin=1in]{geometry}
\usepackage{amsmath, amssymb, mathtools}
\usepackage{graphicx}
\usepackage{float}
\usepackage{booktabs}
\usepackage{longtable}
\usepackage{enumitem}
\usepackage{siunitx}
\usepackage[dvipsnames]{xcolor}
\usepackage{tcolorbox}
\usepackage{fancyhdr}
\usepackage{caption}
\usepackage{tikz}
\usepackage{pgfplots}
\pgfplotsset{compat=1.18}
\usepackage[colorlinks=true, allcolors=NavyBlue]{hyperref}

\captionsetup{font=small, labelfont=bf}
\setlist[enumerate]{itemsep=2pt, topsep=4pt}
\setlist[itemize]{itemsep=2pt, topsep=4pt}

% --- final-answer box -------------------------------------------------------
% Every problem file ends with its result in one of these. The condensed
% solutions document is assembled by lifting exactly these.
\newtcolorbox{answerbox}[1][]{
  colback=Green!6, colframe=Green!55!black,
  boxrule=0.6pt, arc=2pt, left=6pt, right=6pt, top=4pt, bottom=4pt,
  #1
}

% --- citation to course material --------------------------------------------
% \cite{Module 01, PDF p55}  ->  a small italic parenthetical.
% Solutions cite the SLIDE/PAGE a formula came from so a grader (or a future
% reader) can check the claim against what was actually taught. See the
% directive: every non-obvious step carries one.
\newcommand{\srcref}[1]{{\small\itshape(#1)}}

% --- a step that came from OUTSIDE the course materials ---------------------
% Used only where the course is genuinely silent and the problem asks for
% outside research. Marked visually so it is never mistaken for lecture content.
\newcommand{\extref}[1]{{\small\itshape\color{Sepia}[external: #1]}}

% --- callout for the trap in a problem --------------------------------------
\newtcolorbox{trap}[1][]{
  colback=BurntOrange!7, colframe=BurntOrange!70!black,
  boxrule=0.6pt, arc=2pt, left=6pt, right=6pt, top=4pt, bottom=4pt,
  fonttitle=\bfseries, title=Watch out, #1
}

% --- per-problem header -----------------------------------------------------
% \problemheader{8}{15 pts}{Target clock rate}
% The optional 4-argument form carries a learning outcome where the course
% states one (CESC 410 does; CESC 470's HW1 does not).
\newcommand{\problemheader}[3]{%
  \begin{center}\large\bfseries
    Problem #1 \quad\normalsize\mdseries(#2)\\[2pt]
    \large\bfseries #3
  \end{center}
  \vspace{2pt}\hrule\vspace{8pt}
}
\newcommand{\problemheaderlo}[4]{%
  \begin{center}\large\bfseries
    Problem #1 \quad\normalsize\mdseries(#2, #3)\\[2pt]
    \large\bfseries #4
  \end{center}
  \vspace{2pt}\hrule\vspace{8pt}
}

% Course name in the running head. Defined BEFORE \lhead references it, and
% \renewcommand'd by each course's macros file (inputted after this one).
\providecommand{\coursename}{}

\pagestyle{fancy}
\fancyhf{}
\lhead{\small\coursename}
\rhead{\small\thepage}
\renewcommand{\headrulewidth}{0.3pt}

%   % ─────────────────────────────────────────────────────────────────────────────
% CESC 410 course-specific macros. Inputted AFTER the shared preamble:
%
%   \input{../../../../docs/latex/coursework_preamble.tex}
%   \input{../reference_docs/cesc410_macros.tex}
%   \begin{document}
%
% ONLY things that mean something in DSP go here. The document class, the
% package set, \answerbox, \srcref, \extref, the trap callout, \problemheader
% and the page style are identical across courses and live in the shared
% preamble — do not redefine them here.
%
% Build-check one problem in isolation, or the whole assignment:
%   docs/latex/build_tex.sh content/cesc_410/hw/hw01/p01_phasor_form.tex
%   docs/latex/build_tex.sh content/cesc_410/hw/hw01
% ─────────────────────────────────────────────────────────────────────────────

% Fills the running head defined by the shared preamble.
\renewcommand{\coursename}{CESC 410}

% --- DSP notation -----------------------------------------------------------
% Defined once so every problem writes the same symbols. The course uses a
% cursive w for omega on the board; we render omega throughout.
% DTFT. Takes the SYMBOL so \dtft{Y} gives Y(e^{jw}); previously it declared an
% argument and then ignored it, so \dtft{Y} silently printed X(e^{jw}) -- the
% wrong signal, with no error. Unused in HW1, so nothing rendered was affected.
\newcommand{\dtft}[1]{#1\!\left(e^{j\omega}\right)}
\newcommand{\ztr}[1]{\mathcal{Z}\!\left\{#1\right\}}    % z-transform
\newcommand{\conv}{\ast}                                % convolution
\newcommand{\Real}{\operatorname{Re}}
\newcommand{\Imag}{\operatorname{Im}}
% Magnitude / angle. The \mathord wrapper matters: \angle is a mathord, so a
% following "-" would be parsed as a BINARY minus and typeset as "10 /_ - 75",
% with spacing on both sides. Wrapping starts a fresh math list, making the
% sign unary: "10 /_ -75".
\newcommand{\phasor}[2]{#1\,\angle\,\mathord{#2}}
\newcommand{\ustep}{u}                                  % unit step u[n]
\DeclareMathOperator{\sinc}{sinc}

% --- stem plot helper -------------------------------------------------------
% \stemplot{-2}{5}{{0,1},{1,-0.7},{2,0.5}}   % xmin, xmax, {n,value} pairs
% Matches the plotting format the lecture notes use: open circles on stems,
% zero line drawn through.
\newcommand{\stemplot}[4][]{%
  \begin{tikzpicture}[baseline]
    \begin{axis}[
      width=0.62\textwidth, height=4.4cm,
      xmin=#2-0.6, xmax=#3+0.6,
      axis lines=middle, xlabel={$n$}, ylabel={$x[n]$},
      % Park the y-label at the top-left corner of the axis box. With
      % `axis lines=middle` its default position is mid-height ON the y-axis,
      % where a full-height stem at n=0 runs straight through it.
      ylabel style={at={(current axis.north west)}, anchor=south west},
      xtick={#2,...,#3}, tick label style={font=\scriptsize},
      ymajorgrids=false, enlarge y limits=0.25,
      #1
    ]
      \addplot+[ycomb, mark=o, mark options={fill=white}, thick] coordinates {#4};
    \end{axis}
  \end{tikzpicture}%
}

%   \begin{document}
%
% ONLY things that mean something in DSP go here. The document class, the
% package set, \answerbox, \srcref, \extref, the trap callout, \problemheader
% and the page style are identical across courses and live in the shared
% preamble — do not redefine them here.
%
% Build-check one problem in isolation, or the whole assignment:
%   docs/latex/build_tex.sh content/cesc_410/hw/hw01/p01_phasor_form.tex
%   docs/latex/build_tex.sh content/cesc_410/hw/hw01
% ─────────────────────────────────────────────────────────────────────────────

% Fills the running head defined by the shared preamble.
\renewcommand{\coursename}{CESC 410}

% --- DSP notation -----------------------------------------------------------
% Defined once so every problem writes the same symbols. The course uses a
% cursive w for omega on the board; we render omega throughout.
% DTFT. Takes the SYMBOL so \dtft{Y} gives Y(e^{jw}); previously it declared an
% argument and then ignored it, so \dtft{Y} silently printed X(e^{jw}) -- the
% wrong signal, with no error. Unused in HW1, so nothing rendered was affected.
\newcommand{\dtft}[1]{#1\!\left(e^{j\omega}\right)}
\newcommand{\ztr}[1]{\mathcal{Z}\!\left\{#1\right\}}    % z-transform
\newcommand{\conv}{\ast}                                % convolution
\newcommand{\Real}{\operatorname{Re}}
\newcommand{\Imag}{\operatorname{Im}}
% Magnitude / angle. The \mathord wrapper matters: \angle is a mathord, so a
% following "-" would be parsed as a BINARY minus and typeset as "10 /_ - 75",
% with spacing on both sides. Wrapping starts a fresh math list, making the
% sign unary: "10 /_ -75".
\newcommand{\phasor}[2]{#1\,\angle\,\mathord{#2}}
\newcommand{\ustep}{u}                                  % unit step u[n]
\DeclareMathOperator{\sinc}{sinc}

% --- stem plot helper -------------------------------------------------------
% \stemplot{-2}{5}{{0,1},{1,-0.7},{2,0.5}}   % xmin, xmax, {n,value} pairs
% Matches the plotting format the lecture notes use: open circles on stems,
% zero line drawn through.
\newcommand{\stemplot}[4][]{%
  \begin{tikzpicture}[baseline]
    \begin{axis}[
      width=0.62\textwidth, height=4.4cm,
      xmin=#2-0.6, xmax=#3+0.6,
      axis lines=middle, xlabel={$n$}, ylabel={$x[n]$},
      % Park the y-label at the top-left corner of the axis box. With
      % `axis lines=middle` its default position is mid-height ON the y-axis,
      % where a full-height stem at n=0 runs straight through it.
      ylabel style={at={(current axis.north west)}, anchor=south west},
      xtick={#2,...,#3}, tick label style={font=\scriptsize},
      ymajorgrids=false, enlarge y limits=0.25,
      #1
    ]
      \addplot+[ycomb, mark=o, mark options={fill=white}, thick] coordinates {#4};
    \end{axis}
  \end{tikzpicture}%
}

%   \begin{document}
%
% ONLY things that mean something in DSP go here. The document class, the
% package set, \answerbox, \srcref, \extref, the trap callout, \problemheader
% and the page style are identical across courses and live in the shared
% preamble — do not redefine them here.
%
% Build-check one problem in isolation, or the whole assignment:
%   docs/latex/build_tex.sh content/cesc_410/hw/hw01/p01_phasor_form.tex
%   docs/latex/build_tex.sh content/cesc_410/hw/hw01
% ─────────────────────────────────────────────────────────────────────────────

% Fills the running head defined by the shared preamble.
\renewcommand{\coursename}{CESC 410}

% --- DSP notation -----------------------------------------------------------
% Defined once so every problem writes the same symbols. The course uses a
% cursive w for omega on the board; we render omega throughout.
% DTFT. Takes the SYMBOL so \dtft{Y} gives Y(e^{jw}); previously it declared an
% argument and then ignored it, so \dtft{Y} silently printed X(e^{jw}) -- the
% wrong signal, with no error. Unused in HW1, so nothing rendered was affected.
\newcommand{\dtft}[1]{#1\!\left(e^{j\omega}\right)}
\newcommand{\ztr}[1]{\mathcal{Z}\!\left\{#1\right\}}    % z-transform
\newcommand{\conv}{\ast}                                % convolution
\newcommand{\Real}{\operatorname{Re}}
\newcommand{\Imag}{\operatorname{Im}}
% Magnitude / angle. The \mathord wrapper matters: \angle is a mathord, so a
% following "-" would be parsed as a BINARY minus and typeset as "10 /_ - 75",
% with spacing on both sides. Wrapping starts a fresh math list, making the
% sign unary: "10 /_ -75".
\newcommand{\phasor}[2]{#1\,\angle\,\mathord{#2}}
\newcommand{\ustep}{u}                                  % unit step u[n]
\DeclareMathOperator{\sinc}{sinc}

% --- stem plot helper -------------------------------------------------------
% \stemplot{-2}{5}{{0,1},{1,-0.7},{2,0.5}}   % xmin, xmax, {n,value} pairs
% Matches the plotting format the lecture notes use: open circles on stems,
% zero line drawn through.
\newcommand{\stemplot}[4][]{%
  \begin{tikzpicture}[baseline]
    \begin{axis}[
      width=0.62\textwidth, height=4.4cm,
      xmin=#2-0.6, xmax=#3+0.6,
      axis lines=middle, xlabel={$n$}, ylabel={$x[n]$},
      % Park the y-label at the top-left corner of the axis box. With
      % `axis lines=middle` its default position is mid-height ON the y-axis,
      % where a full-height stem at n=0 runs straight through it.
      ylabel style={at={(current axis.north west)}, anchor=south west},
      xtick={#2,...,#3}, tick label style={font=\scriptsize},
      ymajorgrids=false, enlarge y limits=0.25,
      #1
    ]
      \addplot+[ycomb, mark=o, mark options={fill=white}, thick] coordinates {#4};
    \end{axis}
  \end{tikzpicture}%
}


\begin{document}

\problemheaderlo{5}{LO04}{40 pts total, 2 pts each}{Linearity, time invariance, causality, BIBO stability}

\subsection*{Problem statement}

Determine if each of the systems below is (1) linear, (2) time-invariant,
(3) causal, and (4) BIBO stable. \emph{(No procedure is needed if you know how
to perform them.)}
\begin{enumerate}
  \item $T_a(x[n]) = n\,x[n]$
  \item $T_b(x[n]) = a\,x[n] + c$, where $a$ and $c$ are constants
  \item $T_c(x[n]) = x[n+1] + x[n-1]$
  \item $T_d(x[n]) = x[n]\,x[n+1]$
  \item $T_e(x[n]) = e^{x[n-1]}$
\end{enumerate}

\subsection*{Approach}

Four independent tests, applied to each system. The handout says no procedure is
required, but the work is shown here because a single counterexample is what
actually settles each ``no'' — asserting it is not the same as proving it.

\paragraph{Linearity} requires both additivity and homogeneity, tested together
via superposition:
\begin{equation}
  T\{\alpha x_1[n] + \beta x_2[n]\} \overset{?}{=} \alpha T\{x_1[n]\} + \beta T\{x_2[n]\}.
  \label{eq:lin}
\end{equation}
A fast necessary condition: a linear system must map the zero input to the zero
output. If $T\{0\} \ne 0$, it is not linear, no further work needed.

\paragraph{Time invariance} requires that delaying the input delays the output by
the same amount and does nothing else. Compare
\begin{equation}
  y_1[n] = T\{x[n-n_0]\}
  \quad\text{against}\quad
  y_2[n] = y[n-n_0] = \big(T\{x[n]\}\big)\big|_{n \to n-n_0}.
  \label{eq:ti}
\end{equation}
The classic failure is a coefficient in $n$ that does \emph{not} get shifted,
because it belongs to the system rather than to the signal.

\paragraph{Causality} requires $y[n]$ to depend only on $x[m]$ for $m \le n$.
Any appearance of $x[n+1]$ or later is disqualifying.

\paragraph{BIBO stability} requires that every bounded input produce a bounded
output: if $|x[n]| \le B < \infty$ for all $n$, then $|y[n]| \le M < \infty$.

\subsection*{Work}

\paragraph{1. $T_a(x[n]) = n\,x[n]$.}

\emph{Linear — YES.}
\begin{equation}
  T\{\alpha x_1 + \beta x_2\} = n\big(\alpha x_1[n] + \beta x_2[n]\big)
    = \alpha\,\underbrace{n x_1[n]}_{T\{x_1\}} + \beta\,\underbrace{n x_2[n]}_{T\{x_2\}}. \checkmark
\end{equation}

\emph{Time-invariant — NO.} The multiplier $n$ is a property of the system, not
of the signal, so it does not travel with a shift:
\begin{align}
  T\{x[n-n_0]\} &= n\,x[n-n_0], \\
  y[n-n_0]      &= (n-n_0)\,x[n-n_0].
\end{align}
These differ by $n_0 x[n-n_0] \ne 0$. Concretely, with $x[n] = \delta[n]$ and
$n_0 = 1$: $T\{\delta[n-1]\} = n\delta[n-1] = \delta[n-1]$ (value 1 at $n=1$),
while $y[n-1] = (n-1)\delta[n-1] = 0$ everywhere. A nonzero signal versus the
zero signal — not time-invariant.

\emph{Causal — YES.} $y[n]$ uses only $x[n]$, the present sample.

\emph{BIBO stable — NO.} Take the bounded input $x[n] = 1$ for all $n$, so
$B = 1$. Then $y[n] = n$, which grows without bound. The gain $n$ is unbounded
over $n$, so no finite $M$ exists.

\paragraph{2. $T_b(x[n]) = a\,x[n] + c$.}

\emph{Linear — NO} (for $c \ne 0$). This is \emph{affine}, not linear. The zero
input gives $T\{0\} = c \ne 0$, which already fails. Via superposition:
\begin{align}
  T\{\alpha x_1 + \beta x_2\} &= a(\alpha x_1 + \beta x_2) + c, \\
  \alpha T\{x_1\} + \beta T\{x_2\} &= a\alpha x_1 + a\beta x_2 + (\alpha + \beta)c,
\end{align}
which agree only when $(\alpha+\beta)c = c$, i.e.\ only for special $\alpha,\beta$
— not for all. \textbf{If $c = 0$ the system is linear}; the answer below assumes
$c \ne 0$, since otherwise the constant would not have been mentioned.

\emph{Time-invariant — YES.} Both $a$ and $c$ are constants, carrying no
dependence on $n$:
\begin{equation}
  T\{x[n-n_0]\} = a\,x[n-n_0] + c = y[n-n_0]. \checkmark
\end{equation}

\emph{Causal — YES.} Only $x[n]$ is used.

\emph{BIBO stable — YES.} If $|x[n]| \le B$ then by the triangle inequality
\begin{equation}
  |y[n]| \le |a|\,|x[n]| + |c| \le |a|B + |c| < \infty,
\end{equation}
a finite bound independent of $n$.

\paragraph{3. $T_c(x[n]) = x[n+1] + x[n-1]$.}

\emph{Linear — YES.} A sum of shifted inputs with constant coefficients:
\begin{equation}
  T\{\alpha x_1 + \beta x_2\}
    = \alpha x_1[n+1] + \beta x_2[n+1] + \alpha x_1[n-1] + \beta x_2[n-1]
    = \alpha T\{x_1\} + \beta T\{x_2\}. \checkmark
\end{equation}

\emph{Time-invariant — YES.} The shifts are by constants, so
\begin{equation}
  T\{x[n-n_0]\} = x[n-n_0+1] + x[n-n_0-1] = y[n-n_0]. \checkmark
\end{equation}

\emph{Causal — NO.} $y[n]$ depends on $x[n+1]$, one sample into the future. At
$n = 0$ the output needs $x[1]$, which has not arrived yet.

\emph{BIBO stable — YES.} If $|x[n]| \le B$ then
$|y[n]| \le |x[n+1]| + |x[n-1]| \le 2B < \infty$.

\paragraph{4. $T_d(x[n]) = x[n]\,x[n+1]$.}

\emph{Linear — NO.} The input appears multiplied by itself, so scaling the input
by $\alpha$ scales the output by $\alpha^2$:
\begin{equation}
  T\{\alpha x[n]\} = (\alpha x[n])(\alpha x[n+1]) = \alpha^2\,T\{x[n]\}
    \ne \alpha\,T\{x[n]\}
\end{equation}
unless $\alpha \in \{0,1\}$. Homogeneity fails, so the system is nonlinear.

\emph{Time-invariant — YES.} No explicit $n$ appears as a coefficient:
\begin{equation}
  T\{x[n-n_0]\} = x[n-n_0]\,x[n-n_0+1] = y[n-n_0]. \checkmark
\end{equation}
Nonlinear and time-invariant are independent properties — a system can be either
without the other.

\emph{Causal — NO.} $x[n+1]$ is a future value.

\emph{BIBO stable — YES.} If $|x[n]| \le B$ then
$|y[n]| = |x[n]|\,|x[n+1]| \le B^2 < \infty$. A finite bound squared is still
finite.

\paragraph{5. $T_e(x[n]) = e^{x[n-1]}$.}

\emph{Linear — NO.} The exponential is nonlinear. Immediately,
$T\{0\} = e^0 = 1 \ne 0$, which fails the zero-input test. Explicitly,
$e^{x_1 + x_2} = e^{x_1}e^{x_2} \ne e^{x_1} + e^{x_2}$ in general.

\emph{Time-invariant — YES.} The delay is by a constant and the nonlinearity is
memoryless in form:
\begin{equation}
  T\{x[n-n_0]\} = e^{x[n-n_0-1]} = y[n-n_0]. \checkmark
\end{equation}

\emph{Causal — YES.} $y[n]$ depends on $x[n-1]$, a past sample only.

\emph{BIBO stable — YES.} If $|x[n]| \le B$ then $-B \le x[n-1] \le B$, and since
$e^{(\cdot)}$ is monotonically increasing,
\begin{equation}
  0 < e^{-B} \le |y[n]| = e^{x[n-1]} \le e^{B} < \infty .
\end{equation}
The bound $e^B$ is large but finite, which is all stability requires.

\subsection*{Check}

Two structural observations that catch a mis-filled table:

\begin{itemize}
  \item \textbf{Only $T_a$ is unstable.} It is the only system whose gain depends
    on $n$ without bound; every other system has coefficients bounded over $n$,
    so a bounded input cannot escape.
  \item \textbf{Exactly the systems containing $x[n+1]$ are non-causal} — $T_c$
    and $T_d$. Causality is decidable by inspection of the indices alone.
\end{itemize}

Cross-checking linearity against the zero-input test: $T_b\{0\} = c \ne 0$,
$T_e\{0\} = 1 \ne 0$ — both correctly marked nonlinear. $T_a\{0\} = 0$,
$T_c\{0\} = 0$, $T_d\{0\} = 0$, consistent with their entries (necessary, not
sufficient, but no contradiction).

\subsection*{Answer}

\begin{answerbox}
\begin{center}
\begin{tabular}{clcccc}
  \toprule
  & System & Linear & Time-invariant & Causal & BIBO stable \\
  \midrule
  1 & $T_a = n\,x[n]$          & \textbf{Yes} & \textbf{No}  & \textbf{Yes} & \textbf{No}  \\
  2 & $T_b = a\,x[n] + c$      & \textbf{No}$^\dagger$ & \textbf{Yes} & \textbf{Yes} & \textbf{Yes} \\
  3 & $T_c = x[n+1] + x[n-1]$  & \textbf{Yes} & \textbf{Yes} & \textbf{No}  & \textbf{Yes} \\
  4 & $T_d = x[n]\,x[n+1]$     & \textbf{No}  & \textbf{Yes} & \textbf{No}  & \textbf{Yes} \\
  5 & $T_e = e^{x[n-1]}$       & \textbf{No}  & \textbf{Yes} & \textbf{Yes} & \textbf{Yes} \\
  \bottomrule
\end{tabular}
\end{center}
\vspace{2pt}
{\footnotesize $^\dagger$ Affine, not linear, for $c \ne 0$. If $c = 0$ it is linear.}
\end{answerbox}

\end{document}
